% TOP_PROBLEM: {"id": 30006973, "title": "A certified bound of 236 for bounded prime gaps", "category_id": 1, "category": {"id": 1, "name": "number_theory", "display_name": "Number Theory", "description": "Properties of integers, prime numbers, Diophantine equations.", "slug": "number-theory", "order_index": 1, "created_at": "2026-07-31T15:26:25.670Z"}, "status": "open"}
% FIRST_POSTED: 2026-09-27
% ENTRY_KIND: statement_only
% Research continuation by GPT-5.6 Pro, 27 September 2026.
% The established source theorem H_1 <= 240 is preserved; the k=48 / 236 work below is only a numerical candidate until an exact or interval-arithmetic certificate is supplied.
\documentclass[11pt]{article}
\usepackage[margin=25mm]{geometry}
\usepackage{fontspec}
\IfFontExistsTF{Latin Modern Roman}{\setmainfont{Latin Modern Roman}}{\setmainfont{DejaVu Serif}}
\usepackage{amsmath,amssymb,mathtools}
\usepackage{unicode-math}
\IfFontExistsTF{Latin Modern Math}{\setmathfont{Latin Modern Math}}{\setmathfont{DejaVu Math TeX Gyre}}
\usepackage{xurl,hyperref}
\hypersetup{colorlinks=true,urlcolor=blue}
\setcounter{secnumdepth}{0}
\setlength{\parindent}{0pt}
\setlength{\parskip}{5pt}
\setlength{\emergencystretch}{3em}
\begin{document}
% problemId: 30006973
% source status: Julia Stadlmann, arXiv:2608.31126v1, proves H_1 <= 240.
\section*{\texorpdfstring{A certified bound of 236 for bounded prime gaps}{A certified bound of 236 for bounded prime gaps}}\label{prime-gaps-below-240}\label{30006973}

\subsection{Definitions and mathematical statement}
Let $p_n$ denote the $n$th prime and put
\[
 H_1=\liminf_{n\to\infty}(p_{n+1}-p_n).
\]
An increasing $k$-tuple $\mathcal H=(h_1,\ldots,h_k)$ of integers is \emph{admissible} if, for every prime $p$, the residues $h_1,\ldots,h_k\pmod p$ do not occupy every residue class.  Write $H(k)$ for the minimum possible diameter $h_k-h_1$ of an admissible $k$-tuple.

Julia Stadlmann's 2026 preprint proves
\[
 \boxed{H_1\le 240},
\]
by combining Bombieri--Vinogradov with newer equidistribution estimates for suitable moduli and then verifying the GPY/Maynard--Tao optimization for $k=49$ \href{https://arxiv.org/abs/2608.31126}{[S1]}.  The research target of this continuation is to push the unconditional bound strictly below $240$.

The concrete candidate investigated in this session is
\[
 \boxed{H_1\le 236},
\]
via a $k=48$ sieve calculation.  The exact target is therefore to produce a rigorously certified admissible $48$-tuple of diameter at most $236$ together with a rigorously certified sieve weight satisfying the required positivity inequality.  The numerical work described below does \emph{not} yet establish this assertion.

\par\smallskip\textit{Formulation and concept references:} \href{https://arxiv.org/abs/2608.31126}{[S1]}, \href{https://arxiv.org/abs/1407.4897}{[S2]}.

\subsection{Short English statement}
Can the newly proved bound saying that infinitely many consecutive prime gaps are at most $240$ be improved to at most $236$ by reducing the sieve dimension from $k=49$ to $k=48$ and rigorously certifying a stronger symmetric sieve weight?

\subsection{Sources}
\begin{itemize}
\item[S1] Julia Stadlmann, \emph{Bounded gaps between primes}, arXiv:2608.31126v1, submitted 31 August 2026.
\url{https://arxiv.org/abs/2608.31126}.
\textit{Primary source for the proved bound $H_1\le240$, the enlarged-support GPY criterion, the $k=49$ matrix computation, and the final support parameters. Consulted 27 September 2026.}

\item[S2] D. H. J. Polymath, \emph{Variants of the Selberg sieve, and bounded intervals containing many primes}, Research in the Mathematical Sciences 1 (2014), 1--83; arXiv:1407.4897.
\url{https://arxiv.org/abs/1407.4897}.
\textit{Background source for the Maynard--Tao/Polymath optimization and the earlier $H_1\le246$ result.}

\item[S3] Andrew V. Sutherland and contributors, \emph{Narrow admissible tuples} database.
\url{https://math.mit.edu/~primegaps/}.
\textit{Public computational reference for narrow admissible tuples.  The explicit $k=48$ tuple used for the final certificate should be copied into the proof and independently checked rather than relying only on a database lookup. Consulted 27 September 2026.}

\item[E1] Polymath8b supplementary reference material and archived numerical optimization notebooks, hosted in Terence Tao's public repository.
\url{https://github.com/teorth/tao-web/tree/master/static/papers/polymath8b}.
\textit{Exploratory cross-check: the historical optimizer targets the critical variational threshold in dimension $k=50$. Consulted 27 September 2026.}

\item[E2] Axiom Math, \emph{PrimeGapsLib}, public Lean library for formalizations related to prime gaps.
\url{https://github.com/AxiomMath/PrimeGapsLib}.
\textit{Exploratory formalization reference.  The repository records a formal $H_1\le246$ result and a separation between the theoretical development and large numerical certificates; it does not certify the present $236$ candidate. Consulted 27 September 2026.}

\item[N1] Numerical $k=48$ experiment from the present research session, 27 September 2026.
\textit{Unpublished exploratory computation.  Repeated floating-point/grid runs suggested a variational excess of order $10^{-3}$, but no fixed exact coefficient vector, interval certificate, or complete error budget was preserved in the transcript.  It is evidence for a candidate only, not a source for a theorem.}
\end{itemize}

\end{document}
